%
%Academic CV LaTeX Template
% Author: Dubasi Pavan Kumar
% LinkedIn: https://www.linkedin.com/in/im-pavankumar/
% License: MIT
%
% For errors, suggestions, or improvements, please contact:
% Email: pavankumard.pg19.ma@nitp.ac.in
%



\documentclass[a4paper,11pt]{article}

% Package imports
\usepackage{latexsym}
\usepackage{xcolor}
\usepackage{float}
\usepackage{ragged2e}
\usepackage[empty]{fullpage}
\usepackage{wrapfig}
\usepackage{lipsum}
\usepackage{tabularx}
\usepackage{titlesec}
\usepackage{geometry}
\usepackage{marvosym}
\usepackage{verbatim}
\usepackage{enumitem}
\usepackage{fancyhdr}
\usepackage{multicol}
\usepackage{graphicx}
\usepackage{cfr-lm}
\usepackage[T1]{fontenc}
\usepackage{fontawesome5}
\usepackage{multicol}


% Color definitions
\definecolor{darkblue}{RGB}{0,0,139}

% Page layout
\setlength{\multicolsep}{0pt} 
\pagestyle{fancy}
\fancyhf{} % clear all header and footer fields
\fancyfoot{}
\renewcommand{\headrulewidth}{0pt}
\renewcommand{\footrulewidth}{0pt}
\geometry{left=1.4cm, top=0.8cm, right=1.2cm, bottom=1cm}
\setlength{\footskip}{5pt} % Addressing fancyhdr warning

% Hyperlink setup (moved after fancyhdr to address warning)
\usepackage[hidelinks]{hyperref}
\hypersetup{
    colorlinks=true,
    linkcolor=darkblue,
    filecolor=darkblue,
    urlcolor=darkblue,
}

% Custom box settings
\usepackage[most]{tcolorbox}
\tcbset{
    frame code={},
    center title,
    left=0pt,
    right=0pt,
    top=0pt,
    bottom=0pt,
    colback=gray!20,
    colframe=white,
    width=\dimexpr\textwidth\relax,
    enlarge left by=-2mm,
    boxsep=4pt,
    arc=0pt,outer arc=0pt,
}

% URL style
\urlstyle{same}

% Text alignment
\raggedright
\setlength{\tabcolsep}{0in}

% Section formatting
\titleformat{\section}{
  \vspace{-4pt}\scshape\raggedright\large
}{}{0em}{}[\color{black}\titlerule \vspace{-7pt}]

% Custom commands
\newcommand{\resumeItem}[2]{
  \item{
    \textbf{#1}{\hspace{0.5mm}#2 \vspace{-0.5mm}}
  }
}

\newcommand{\resumePOR}[3]{
\vspace{0.5mm}\item
    \begin{tabular*}{0.97\textwidth}[t]{l@{\extracolsep{\fill}}r}
        \textbf{#1}\hspace{0.3mm}#2 & \textit{\small{#3}} 
    \end{tabular*}
    \vspace{-2mm}
}

\newcommand{\resumeSubheading}[4]{
\vspace{0.5mm}\item
    \begin{tabular*}{0.98\textwidth}[t]{l@{\extracolsep{\fill}}r}
        \textbf{#1} & \textit{\footnotesize{#4}} \\
        \textit{\footnotesize{#3}} &  \footnotesize{#2}\\
    \end{tabular*}
    \vspace{-2.4mm}
}

\newcommand{\resumeProject}[4]{
\vspace{0.5mm}\item
    \begin{tabular*}{0.98\textwidth}[t]{l@{\extracolsep{\fill}}r}
        \textbf{#1} & \textit{\footnotesize{#3}} \\
        \footnotesize{\textit{#2}} & \footnotesize{#4}
    \end{tabular*}
    \vspace{-2.4mm}
}

\newcommand{\resumeSubItem}[2]{\resumeItem{#1}{#2}\vspace{-4pt}}

\renewcommand{\labelitemi}{$\vcenter{\hbox{\tiny$\bullet$}}$}
\renewcommand{\labelitemii}{$\vcenter{\hbox{\tiny$\circ$}}$}

\newcommand{\resumeSubHeadingListStart}{\begin{itemize}[leftmargin=*,labelsep=1mm]}
\newcommand{\resumeHeadingSkillStart}{\begin{itemize}[leftmargin=*,itemsep=1.7mm, rightmargin=2ex]}
\newcommand{\resumeItemListStart}{\begin{itemize}[leftmargin=*,labelsep=1mm,itemsep=0.5mm]}

\newcommand{\resumeSubHeadingListEnd}{\end{itemize}\vspace{2mm}}
\newcommand{\resumeHeadingSkillEnd}{\end{itemize}\vspace{-2mm}}
\newcommand{\resumeItemListEnd}{\end{itemize}\vspace{-2mm}}
\newcommand{\cvsection}[1]{%
\vspace{2mm}
\begin{tcolorbox}
    \textbf{\large #1}
\end{tcolorbox}
    \vspace{-4mm}
}

\newcolumntype{L}{>{\raggedright\arraybackslash}X}%
\newcolumntype{R}{>{\raggedleft\arraybackslash}X}%
\newcolumntype{C}{>{\centering\arraybackslash}X}%

% Commands for icon sizing and positioning
\newcommand{\socialicon}[1]{\raisebox{-0.05em}{\resizebox{!}{1em}{#1}}}
\newcommand{\ieeeicon}[1]{\raisebox{-0.3em}{\resizebox{!}{1.3em}{#1}}}

% Font options
\newcommand{\headerfonti}{\fontfamily{phv}\selectfont} % Helvetica-like (similar to Arial/Calibri)
\newcommand{\headerfontii}{\fontfamily{ptm}\selectfont} % Times-like (similar to Times New Roman)
\newcommand{\headerfontiii}{\fontfamily{ppl}\selectfont} % Palatino (elegant serif)
\newcommand{\headerfontiv}{\fontfamily{pbk}\selectfont} % Bookman (readable serif)
\newcommand{\headerfontv}{\fontfamily{pag}\selectfont} % Avant Garde-like (similar to Trebuchet MS)
\newcommand{\headerfontvi}{\fontfamily{cmss}\selectfont} % Computer Modern Sans Serif
\newcommand{\headerfontvii}{\fontfamily{qhv}\selectfont} % Quasi-Helvetica (another Arial/Calibri alternative)
\newcommand{\headerfontviii}{\fontfamily{qpl}\selectfont} % Quasi-Palatino (another elegant serif option)
\newcommand{\headerfontix}{\fontfamily{qtm}\selectfont} % Quasi-Times (another Times New Roman alternative)
\newcommand{\headerfontx}{\fontfamily{bch}\selectfont} % Charter (clean serif font)

\begin{document}
\headerfontiii

% Header
\begin{center}
    {\Large\textbf{Duc Hoang}}
\end{center}
\vspace{-6mm}

\begin{center}
    \small{
    +1-901-628-8339 | \href{mailto:dhoang@mit.edu}{dhoang@mit.edu} | 
    \href{https://www.yourwebsite.com/}{Google Scholar}|
    \socialicon{\faGithub} \href{https://github.com/your-username}{Duchstf} 
    | Citizenship: Vietnam
    }
\end{center}
\vspace{-6mm}

\section{\textbf{Summary}}
\vspace{1mm}
\small{
High-energy physicist specializing in machine learning, data science, and real-time FPGA systems. Extensive experience delivering production-level ultra-low-latency ML inference on FPGAs and petabyte-scale data analysis for CERN’s CMS experiment. Proven leader in multi-institutional projects, skilled in developing custom ML solutions and deploying robust, high-performance software.
}
\vspace{-2mm}

\section{\textbf{Education}}
\vspace{-0.4mm}
\resumeSubHeadingListStart

\resumeSubheading
{Massachusetts Institute of Technology (MIT)}{Cambridge, MA}
{PhD, High Energy Physics}{2021 - 2026}

\resumeSubheading
{Rhodes College}{Memphis, USA}
{B.S., Physics, Summa Cum Laude}{2017-2021}

\resumeSubHeadingListEnd
\vspace{-6mm}




\section{\textbf{Professional Experience}}
\vspace{-0.4mm}
  \resumeSubHeadingListStart
  \resumeSubheading
      {{MIT \& European Organization for Nuclear Research (CERN)}}{Cambridge, USA \& Geneva, Switzerland}
      {Graduate Researcher \& Next Generation Trigger PhD Fellow}{2021 - Now}
      \resumeItemListStart
        \item Coordinated multi-institutional team to develop \textbf{ultra-low-latency FPGA-accelerated ML models} for CMS Level-1 Trigger, which processes approximately 10\% of global internet traffic in real time. \href{https://www.youtube.com/watch?v=JvdCE5woYuY}{Hear me explain it along the trigger leardership team} \textbf{[1]}
        \item Built first \textbf{FPGA-based ML algorithms for real-time jet tagging} and delivered cross-platform, low-latency data transfer tests between different FPGA board platforms. \textbf{[2]}
        \item Analyzed all \textbf{Run 2 CMS data (a few PBs)} to deliver \textbf{most sensitive Higgs $\to b\overline{b}$ measurement in the V(qq)H(bb) decay channel}, invited presentation at Rencontres de Moriond 2025. \textbf{[3]}
        \item Created and taught data analysis projects for \textbf{MIT 8.16/8.316: Data Science in Physics}, which is published on \textbf{MITx} course \textbf{\href{https://mitxonline.mit.edu/courses/course-v1:MITxT+8.S50.2x/}{Computational Data Science in Physics}}.
        \item Supervised MIT undergraduates on research projects involving the CMS Level-1 Trigger system and the development of novel machine learning architectures on FPGAs. \textbf{[9]}.
      \resumeItemListEnd 
  \resumeSubheading
    {Fermi National Accelerator Laboratory}{Chicago, USA}
    {Undergraduate Researcher}{2019 - 2021}
    \resumeItemListStart
      \item Developer of  \href{https://github.com/fastmachinelearning/hls4ml}{hls4ml}, A package for \textbf{ultra-low-latency machine learning inference in FPGAs} using high level synthesis language (HLS). \textbf{[4]}
      \item Designed superconducting magnet quench detection systems \textbf{[5]}
      \item Analyzed simulation data to predict potential signal gains for the real-time trigger systems for Light Dark Matter Experiment (LDMX). \textbf{[6]}

      \item Machine Learning analysis to infer hyper-parameter optimization performance of models used in the MINERvA Experiment. \textbf{[7]}
    \resumeItemListEnd
  \resumeSubHeadingListEnd
\vspace{-6mm}

\section*{\textbf{Selected Publications \& Software}}
\vspace{0.2mm}
\small{
\begin{multicols}{2}
\begin{enumerate}[leftmargin=*, labelsep=0.5em, align=left, widest={[\textbf{S.1}]}, itemindent=0em, label={\textbf{[\arabic*]}}]

\item \textbf{CMS Collaboration}, ``The Phase-2 Upgrade of the CMS Level-1 Trigger", \url{https://cds.cern.ch/record/2714892}

\item \textbf{CMS Collaboration}, 2025. ``Multiclass object tagging in the CMS Phase-2 Level-1 Correlator Trigger." CMS Detector Performance Summary. \url{https://cds.cern.ch/record/2936315.}

\item \textbf{CMS Collaboration}, ``Search for a boosted Higgs boson decaying to bottom quark pairs in association with a hadronically decaying W or Z boson with the CMS detector using proton-proton collisions at $\sqrt{s}=13$ TeV", \url{https://cds.cern.ch/record/2928061}

\item \texttt{hls4ml} \textbf{Community}, \href{[arXiv:2103.05579 [cs.LG]].}{``hls4ml: An Open-Source Codesign Workflow to Empower Scientific Low-Power Machine Learning Devices"}


\item \textbf{D. Hoang} et al., ``Intelliquench: an adaptive machine learning system for detection of superconducting magnet quenches.," in \emph{IEEE Transactions on Applied Superconductivity}, doi: \url{10.1109/TASC.2021.3058229}.

\item \textbf{LDMX Collaboration}, ``Current Status and Future Prospects for the Light Dark Matter eXperiment" \url{[arXiv:2203.08192 [hep-ex]]}.

\item \textbf{D. Hoang} et al. ``Inferring Convolutional Neural Networks' Accuracies from Their Architectural Characterizations," \emph{2019 18th IEEE International Conference On Machine Learning And Applications (ICMLA)}, Boca Raton, FL, USA, 2019, pp. 1388-1391. DOI: \url{10.1109/ICMLA.2019.00226}

\item \textbf{Dark Quest Collaboration}, ``DarkQuest: A dark sector upgrade to SpinQuest at the 120 GeV Fermilab Main Injector" \url{[arXiv:2203.08322 [hep-ex]]}.

\item \textbf{CMS Collaboration}. 2024. ``NNPuppiTaus: PUPPI Tau Reconstruction in the Level-1 Trigger with Real-Time Machine Learning." CMS Detector Performance Summary. \url{https://cds.cern.ch/record/2895660}.

\end{enumerate}
\end{multicols}



\section{\textbf{Skills}}
\vspace{-0.4mm}
 \resumeHeadingSkillStart
    \resumeSubItem{Specialized Area:}
    {Big Data Analysis, Statistical Modeling, ultra-low latency FPGA Inference, Machine Learning}
  \resumeSubItem{Programming Languages:}
    {Python, C++, HTML, CSS, TypeScript}
  \resumeSubItem{Web Technologies:}
    {Django, Angular}
  \resumeSubItem{Database Systems:}
    {SQLite, PostgreSQL}
  \resumeSubItem{Data Science \& Machine Learning:}
    {Most Python-related machine learning tools.}
\resumeSubItem{Cloud Technologies:}
    {Google Cloud}
  \resumeSubItem{DevOps \& Version Control:}
    {Continous Integration, Docker, Singularity/Apptainer, Git}
 \resumeHeadingSkillEnd

\section{\textbf{Honors and Awards}}
\vspace{-0.4mm}
\resumeSubHeadingListStart

\resumeProject
  {The 2025 Breakthrough Prize in Fundamental Physics}
  {Received as a member of the CMS Collaboration}
  {2025}
  {{}[\href{https://breakthroughprize.org/Laureates/1/L3994}{\textcolor{darkblue}{\faIcon{globe}}}]}
\resumeItemListStart
  \item The prize is awarded to physicists from theoretical, mathematical, or experimental physics that have made transformative contributions to fundamental physics, and specifically for recent advances. 
\resumeItemListEnd

\resumeProject
  {Next Generation Trigger PhD fellow.}
  {CERN}
  {2024}
  {{}[\href{https://cerncourier.com/a/next-generation-triggers-for-hl-lhc-and-beyond/}{\textcolor{darkblue}{\faIcon{globe}}}]}
\resumeItemListStart
  \item Awarded for exceptional PhD students to work along CERN scientists to develop the next generation trigger system at the Large Hadron Collider.
\resumeItemListEnd

\resumeProject
  {MIT Public Service Fellow (PKG) Fellowship}
  {Priscilla King Gray Public Service Center - MIT}
  {2025}
  {{}[\href{https://pkgcenter.mit.edu/programs/fellowships/pkg-fellowships/}{\textcolor{darkblue}{\faIcon{globe}}}]}
\resumeItemListStart
  \item Founded and organized, along with other International Olympiad Winners from Central Vietnam, a STEM \href{https://seas-cvn.com/}{summer school}, providing under-privileged high school students in Central Vietnam the opportunity to study emerging fields such as artificial intelligence, data science, programming, renewable energy, and quantum technology at no cost.
\resumeItemListEnd

\resumeProject
  {2024-25 Christiaan Huygens Graduate Fellowship - MIT.}
  {MIT School of Science}
  {2025}
  {}
\resumeItemListStart
  \item Internal Fellowship providing financial support for the full academic year in the School of Science for students with non-traditional backgrounds.
\resumeItemListEnd

\resumeProject
  {MIT-The Royal Game of Ur Champion}
  {MIT Playur II Tournament}
  {2023}
  {{}[\href{https://www.instagram.com/p/CsHxpyhuZw3/?img_index=1}{\textcolor{darkblue}{\faIcon{globe}}}]}
\resumeItemListStart
  \item The Royal Game of Ur is the oldest board game in the world. 
\resumeItemListEnd

\resumeProject
  {Henry W Kendall (1955) first year graduate fellowship.}
  {MIT Laboratory for Nuclear Science}
  {2021}
  {}
\resumeItemListStart
  \item Financial support for first year graduate students.
\resumeItemListEnd

\resumeProject
  {Miscellaneous Undergraduate Awards.}
  {Rhodes College}
  {2017-2021}
  {}
\resumeItemListStart
  \item  Research Award in Physics, Council on Undergraduate Research (CUR) Physics \& Astronomy Student Travel Award, W.O Shewmaker Award for Excellence in First-Year Search Program,  Award for Excellence in Physics by a First-Year Student, Summa Cum Laude, Sigma Pi Sigma, Phi Beta Kappa, Physics Honors.
\resumeItemListEnd

\resumeSubHeadingListEnd

\end{document}